\documentclass[11pt,a4paper]{article}

\usepackage[utf8]{inputenc}
\usepackage[T1]{fontenc}
\usepackage{amsmath,amssymb,amsthm}
\usepackage{mathtools}
\usepackage{booktabs}
\usepackage{array}
\usepackage{geometry}
\usepackage{hyperref}
\usepackage{cleveref}
\usepackage{enumitem}
\usepackage{caption}
\usepackage{xcolor}
\usepackage{microtype}

\geometry{margin=1in}

\newtheorem{theorem}{Theorem}[section]
\newtheorem{lemma}[theorem]{Lemma}
\newtheorem{proposition}[theorem]{Proposition}
\newtheorem{corollary}[theorem]{Corollary}
\newtheorem{definition}[theorem]{Definition}
\newtheorem{remark}[theorem]{Remark}
\newtheorem{hypothesis}[theorem]{Hypothesis}
\newtheorem{prediction}[theorem]{Prediction}

\newcommand{\Jcost}{J}
\newcommand{\ph}{\varphi}
\newcommand{\tick}{\tau_0}
\newcommand{\tbio}{\tau_{\mathrm{bio}}}
\newcommand{\Ecoh}{E_{\mathrm{coh}}}
\newcommand{\Qsix}{\ensuremath{Q_6}}
\newcommand{\RR}{\mathbb{R}}
\newcommand{\NN}{\mathbb{N}}
\newcommand{\ZZ}{\mathbb{Z}}
\newcommand{\CC}{\mathbb{C}}

\title{\textbf{Toward a Zero-Junk Genome:\\[4pt]
Recognition Science Predictions for DNA\\[4pt]
as Z-Invariant Storage}}

\author{Jonathan Washburn\\
\small Recognition Science Research Institute, Austin, Texas\\
\small \texttt{washburn.jonathan@gmail.com}}

\date{March 17, 2026 (Third-Pass Revision)}

\begin{document}
\maketitle

\begin{abstract}
The genetic code maps $64 = 8^2$ codons to $20$ amino acids, wobble degeneracy dominates $14$ of $16$ codon family boxes, and the $4$\,bp half-octave helical rotation lies within $1^\circ$ of the small golden angle $360^\circ/\ph^2$.
Recognition Science (RS) derives these facts from a single algebraic primitive, the Recognition Composition Law, whose unique cost function forces the golden ratio~$\ph$, spatial dimension $D{=}3$, and an 8-tick discrete clock.
We formalize a $6$-bit codon encoding into a qualia space~\Qsix{} and prove that synonymous wobble substitutions need not preserve Z-parity (the topological invariant of the encoding), establishing a hierarchy: amino-acid identity is conserved more strongly than parity at the wobble position.
We show that the abstract $(8,1,8)$ repetition code saturates the Singleton bound and uniquely maximizes minimum distance among length-$8$ codes.
We present five falsifiable predictions---including mod-$8$ non-uniformity in intron lengths and a minimal-intron window near $8\ph^5 \approx 89$\,bp---and propose a five-layer decomposition of non-coding DNA (error correction, phase alignment, ledger balance, modular control, temporal coordination) as a structured research program.
Since the first draft, the Lean surface has expanded substantially: 49 modules, 8,246 lines, zero axioms, zero unproved assertions.
A major new computation-assisted result---the standard genetic code is optimal over a 363,004-element classwise degeneracy-preserving search under a position-weighted mutation cost---has been added and is described in \S\ref{sec:formal}.
Throughout, we separate formal theorems (machine-checked in Lean~4) from model-level interpretations and empirical hypotheses.
\end{abstract}

\tableofcontents

%============================================================
\section{Introduction}
\label{sec:intro}
%============================================================

Approximately 98\% of the human genome does not encode proteins. What this non-coding majority \emph{does} is the subject of a debate that has run for half a century. Ohno's original ``junk DNA'' proposal~\cite{Ohno1972} treated most non-coding sequence as evolutionary debris. The ENCODE project reported biochemical activity across roughly 80\% of the genome~\cite{ENCODE2012}, but critics argued that detectable activity is not the same as selected function~\cite{Graur2013, Doolittle2013}. Comparative-genomics estimates of the fraction under purifying selection range from 5--15\%~\cite{Ponting2011, Rands2014, Kellis2014}, leaving the status of the remaining non-coding sequence genuinely open.

Recognition Science (RS) enters this debate with an unusual structural claim: the genome is not primarily a protein recipe book but a storage format for topological invariants. In this view, the 8-fold structure of the double helix ($4$ nucleotides $\times$ $2$ strands), the $64 = 8^2$ codon space, the degeneracy pattern, and the non-coding architecture are all reflections of a single algebraic backbone---the Recognition Composition Law (RCL)---from which the golden ratio~$\ph$, spatial dimension $D=3$, and an 8-tick discrete clock are forced consequences~\cite{Washburn2025}.

This paper has three goals:
\begin{enumerate}[nosep]
\item Present the strongest DNA-related results that are currently machine-checked in Lean~4.
\item Propose a five-layer research decomposition of non-coding DNA and state it as a falsifiable program rather than a closed theorem.
\item Identify the concrete empirical and mechanistic steps that separate the present formal scaffold from a complete theory.
\end{enumerate}

We adopt a three-level claim discipline throughout. A \textbf{theorem} is a statement proved in the current Lean development. A \textbf{model} is a structured interpretation implemented as definitions and consistency lemmas. A \textbf{hypothesis} is an empirical or mechanistic claim not yet derived from the formal surface.

%============================================================
\section{Recognition Science Background}
\label{sec:background}
%============================================================

The core primitive of RS is the Recognition Composition Law:
\begin{equation}
\label{eq:rcl}
\Jcost(xy) + \Jcost(x/y) = 2\,\Jcost(x)\,\Jcost(y) + 2\,\Jcost(x) + 2\,\Jcost(y).
\end{equation}
This is a calibrated multiplicative form of the d'Alembert functional equation. Under normalization $\Jcost(1)=0$ and the calibration $\Jcost''(1)=1$, the unique continuous solution is the canonical cost
\begin{equation}
\label{eq:jcanon}
\Jcost(x) = \tfrac{1}{2}(x + x^{-1}) - 1.
\end{equation}
The function $\Jcost$ is non-negative, vanishes only at $x=1$, and satisfies the symmetry $\Jcost(x)=\Jcost(1/x)$.

From this single primitive, a forcing chain (labeled T0--T8 in the Lean codebase) derives a sequence of structural consequences. The steps most relevant to DNA are:
\begin{itemize}[nosep]
\item \textbf{T6} ($\ph$ forcing): Self-similarity in the discrete ledger forces $x^2 = x+1$, so $\ph = (1+\sqrt{5})/2$ appears as the unique positive fixed point. This is the origin of every golden-ratio signature in the theory.
\item \textbf{T7} (8-tick): The minimal ledger-compatible walk in $D=3$ has period $2^3 = 8$, creating an 8-phase discrete clock.
\item \textbf{T8} ($D=3$): Linking requirements and gap-$45$ synchronization ($\operatorname{lcm}(8,45)=360$) force exactly three spatial dimensions.
\end{itemize}

For the biology papers, the practical upshot is that $\ph$, the number~$8$, and three-dimensionality are not free parameters but derived constants. The DNA question becomes: given this algebraic vocabulary, what can be said about the structure of genomes?

%============================================================
\section{Formal Results}
\label{sec:formal}
%============================================================

This section presents the DNA-related statements that are currently machine-checked in Lean~4 within the \texttt{IndisputableMonolith.Genetics} namespace, drawing on the RS constants and combinatorics from the \texttt{RecognitionScience} namespace.

\subsection{Codon Combinatorics}

The double helix provides $4 \times 2 = 8$ recognition slots per position (four nucleotide types on two complementary strands), matching the 8-tick period. The codon space has $4^3 = 64 = 8^2$ entries. Both equalities are proved by computation in Lean.

The standard genetic code maps these 64 codons to 20 amino acids plus 3 stop signals, with degeneracies summing to 61 sense codons. In the RS formalization, 20 ``WTokens'' (semantic atoms derived from the 8-tick DFT basis) are placed in bijection with the 20 canonical amino acids~\cite{Washburn2025}. The cardinality match $20 = 20$ is proved; the specific biochemical assignment (which token maps to which amino acid) is a model-level correspondence, not a derived theorem.

\subsection{Wobble Symmetry}

Third-position degeneracy dominates the code.

\begin{theorem}[14 of 16 Family Boxes]
\label{thm:14of16}
Of the $16$ family boxes defined by the first two codon positions, exactly~$14$ preserve clean wobble symmetry (all four third-position variants synonymous, or a clean pyrimidine/purine $2{+}2$ split).  Only the Ile/Met box ($3{+}1$) and the Cys/Trp/Stop box ($2{+}1{+}1$) break this pattern.  Of the~$14$ clean boxes, exactly~$8$ are fully four-fold degenerate.
\end{theorem}

\noindent
The proof enumerates all $16$ boxes, evaluates the genetic-code mapping at each, and checks the wobble condition by \texttt{native\_decide}.

\begin{theorem}[Wobble Does Not Preserve Z-Parity]
\label{thm:wobble-parity}
In the $6$-bit codon encoding (3 nucleotides $\times$ 2 bits each, mapped into~\Qsix), synonymous wobble substitutions need not preserve the Z-parity (bit-count mod~$2$). Concretely, GTT and GTC both encode valine, but $\mathrm{parity}(\texttt{GTT}) = 1$ while $\mathrm{parity}(\texttt{GTC}) = 0$.
\end{theorem}

This result is interesting because it reveals a hierarchy of conservation in the genetic code. The code prioritizes preserving amino-acid identity over preserving Z-parity at the wobble position. If parity-like invariants play a biological role---as the RS framework proposes---they must operate at a level coarser than single-codon wobble, perhaps across reading frames, exon boundaries, or gene-level sums.

\subsection{The $(8,1,8)$ Repetition Code}

\begin{proposition}
\label{prop:818}
The abstract $(8,1,8)$ repetition code achieves:
\begin{enumerate}[nosep]
\item Maximum possible minimum distance: $d = n = 8$.
\item Tight Singleton bound: $d = n - k + 1$.
\item Error correction of up to $\lfloor(d-1)/2\rfloor = 3$ errors.
\end{enumerate}
Moreover, any $(8,k,d)$ code with $k \geq 2$ has $d \leq 7$, so the repetition code uniquely maximizes distance among length-$8$ codes.
\end{proposition}

This is a coding-theoretic template motivated by the RS 8-tick structure. It establishes that if a biological system uses an 8-phase encoding, the repetition code is the unique distance-optimal architecture. Whether the genetic code literally instantiates this template is an empirical question; the standard code's average degeneracy of $61/20 \approx 3.05$ is consistent with substantial but not maximal redundancy~\cite{HaigHurst1991, FreelandHurst1998}.

\subsection{Helix Geometry and Intron Structure}

\begin{theorem}[Golden-Angle Proximity]
\label{thm:golden-angle}
B-form DNA has approximately $10.5$\,bp per helical turn. The rotation subtended by $4$\,bp (the ``half-octave'') is $4 \times 360^\circ / 10.5 \approx 137.1^\circ$. The small golden angle is $360^\circ / \ph^2 \approx 137.5^\circ$. These differ by less than $1^\circ$:
\[
\left| 4 \times \frac{360}{10.5} - \frac{360}{\ph^2} \right| < 1^\circ.
\]
The proof uses rigorous rational squeezes on $\ph$ (via $\sqrt{5}$ bounds) and is machine-checked.
\end{theorem}

\begin{theorem}[Minimal-Intron Estimate]
\label{thm:min-intron}
Using the algebraic identity $\ph^5 = 5\ph + 3$, the product $8\ph^5$ satisfies $|8\ph^5 - 89| < 1$. The observed minimum functional intron size in mammals is $85$--$90$\,bp.
\end{theorem}

Together, these results support the hypothesis that introns function as helical phase-alignment spacers whose lengths are constrained by the 8-tick clock. The golden-angle proximity explains \emph{why} a 4\,bp unit would be structurally preferred, and the $8\ph^5$ estimate identifies a natural minimal length scale.

\subsection{Global Optimality of the Codon-to-Amino-Acid Assignment}

\begin{proposition}[Computation-Assisted Position-Weighted Classwise Optimality]
\label{thm:global-opt}
Under a position-dependent mutation cost with eight components (hydrophobicity$^2$ at positions 1--3, charge$^2$ at position~3, side-chain volume$^2$ at positions 1--2, and biosynthetic pathway distance at positions 1--2), the standard genetic code is the unique cost minimizer over the 363,004 non-identity permutations obtained by enumerating each degeneracy class separately: $(2!{-}1) + (9!{-}1) + (1!{-}1) + (5!{-}1) + (3!{-}1) = 363{,}004$. The integer-weighted worst-case margin on this classwise search is $+37{,}136$. Lean proves the exact classwise search size and contains 52 \texttt{native\_decide} certificates covering every pairwise swap and every non-identity 6-fold permutation; the full classwise optimality statement is confirmed by exhaustive external computation.
\end{proposition}

This result partially resolves Prediction~\ref{pred:composite} from the first draft: the code \emph{is} locally optimal within every degeneracy class and is exhaustively optimal over the currently enumerated classwise search under a biologically motivated position-weighted metric. The key insight is that mutations at different codon positions have very different chemical consequences: position~2 changes almost always alter the amino acid's chemical class, while position~3 (wobble) changes are frequently synonymous. A companion paper~\cite{Washburn2026b} provides the full derivation, sensitivity analysis, and the exact separation between Lean-certified and externally enumerated parts of the argument.

\subsection{Bio-Clocking and the FOLD Operator}

\begin{proposition}[Bio-Clock Recurrence]
\label{prop:bioclock}
The bio-clocking function $\tbio(N) = \tick \cdot \ph^N$ satisfies $\tbio(N+1) = \ph \cdot \tbio(N)$, and the sequence is strictly monotone increasing.
\end{proposition}

\begin{proposition}[FOLD Idempotence]
\label{prop:fold}
The FOLD operator on the 8-tick register, $\mathrm{FOLD}(v)_k = (v_k + v_{(8-k) \bmod 8})/2$, is idempotent: $\mathrm{FOLD}(\mathrm{FOLD}(v)) = \mathrm{FOLD}(v)$ for all $k \in \{0,\ldots,7\}$. The proof proceeds by exhaustive case analysis over all 8 values.
\end{proposition}

The RS interpretation is that sexual reproduction physically instantiates FOLD: meiotic recombination averages complementary parental phase registers, annihilating accumulated phase drift.  The following two results give this interpretation formal teeth.

\begin{theorem}[FOLD Annihilates Antisymmetric Phase Debt]
\label{thm:fold-cancel}
Decompose any phase profile $v \in \RR^8$ into symmetric and antisymmetric parts:
\[
v^{\mathrm{sym}}_k = \frac{v_k + v_{8-k}}{2},\qquad
v^{\mathrm{anti}}_k = \frac{v_k - v_{8-k}}{2}.
\]
Then $\mathrm{FOLD}(v) = v^{\mathrm{sym}}$ and the antisymmetric part of $\mathrm{FOLD}(v)$ is identically zero for all $k \in \{0,\ldots,7\}$.  The proof proceeds by exhaustive \texttt{fin\_cases} and \texttt{ring}.
\end{theorem}

\begin{corollary}[Sexual Advantage]
\label{cor:sex}
A lineage applying FOLD each generation has bounded (zero) antisymmetric phase debt.  An asexual lineage preserves its antisymmetric component unchanged; if mutations inject positive debt~$\delta$ per generation, the debt grows as $n\delta$ (Muller's ratchet in phase-profile language).
\end{corollary}

\subsection{Position-Specific Mutation Vulnerability}

\begin{theorem}[Wobble is Safest; Position~2 is Most Dangerous]
\label{thm:posorder}
For each of the $64$ codons, compute the total chemical-distance cost of all single-nucleotide mutations at positions~$1$, $2$, and~$3$ separately, weighted by a transition/transversion ratio $\kappa = 2$.  Summing over all codons:
\[
\text{Position 3 (wobble)} \;\leq\; \text{Position 1} \;\leq\; \text{Position 2}.
\]
Both inequalities are strict and are verified by \texttt{native\_decide} over the full $64 \times 9$ mutation table.
\end{theorem}

Position~2 largely determines the chemical class of the amino acid (hydrophobic, polar, charged), so mutations there are the most costly.  The wobble position carries the least information and the most degeneracy---it is the ``least significant bit'' of the codon.  This ordering is a consequence of how the standard code arranges amino acids in the $8 \times 8$ codon grid, and it is one of the properties that any candidate ``J-cost metric'' must reproduce.

%============================================================
\section{A Five-Layer Decomposition of Non-Coding DNA}
\label{sec:layers}
%============================================================

The RS genetics framework suggests a five-layer account of non-coding sequence. We present it here as a \emph{research decomposition}---a structured way to organize existing and future empirical work---rather than as a completed derivation.

\begin{center}
\small
\renewcommand{\arraystretch}{1.3}
\begin{tabular}{@{}p{0.18\linewidth}p{0.40\linewidth}p{0.34\linewidth}@{}}
\toprule
\textbf{Layer} & \textbf{RS interpretation} & \textbf{Current evidence} \\
\midrule
Error correction & Codon degeneracy ($61/20 \approx 3\times$) and the 20/7 token-to-dimension ratio provide noise immunity. & Degeneracy counts formalized; code near-optimality supported by Haig--Hurst and Freeland--Hurst analyses~\cite{HaigHurst1991, FreelandHurst1998}. \\

Phase alignment & Introns act as helical spacers, keeping exon junctions at golden-angle offsets. & Golden-angle theorem (\cref{thm:golden-angle}) and minimal-intron estimate (\cref{thm:min-intron}) are proved; mod-8 enrichment is empirical and dataset-sensitive. \\

Ledger balance & Repetitive elements (SINEs, LINEs, LTR retrotransposons) serve as complementary entries in the $\sigma{=}0$ double-entry ledger. & Hypothesis: conceptually grounded in RS, but not yet tested by targeted knockout or comparative-genomics studies. \\

Modular control & Regulatory DNA is interface architecture, selected by MDL-style fitness $r(g) \propto \exp(-\beta L_g)$. & Model: abstract MDL inequalities proved; genome fractions not derived from annotation data. \\

Temporal coordination & Methylation, replication timing, and FOLD-mediated recombination implement clocks and phasing. & FOLD idempotence and bio-clock recurrence proved; molecular mechanism open. \\
\bottomrule
\end{tabular}
\end{center}

How does this compare with existing accounts of non-coding DNA? The selfish-DNA hypothesis~\cite{Orgel1980} treats transposable elements as genomic parasites; RS reinterprets them as ledger-balancing substrate. The ENCODE ``80\% functional'' claim~\cite{ENCODE2012} counts biochemical activity; RS would require structural necessity under the cost function. The comparative-genomics constraint estimates of 5--15\% under purifying selection~\cite{Ponting2011, Rands2014} measure evolutionary conservation; RS predicts additional \emph{phase-sensitive} constraint that standard alignment-based methods may not detect because it depends on helical position rather than linear sequence identity.

\begin{hypothesis}[Zero-Junk Program]
\label{hyp:zerojunk}
As functional assays, comparative genomics, and phase-sensitive analyses improve, the fraction of sequence assignable to one or more of the five layers will grow substantially beyond current constraint estimates.
\end{hypothesis}

This is not the same as asserting that every base pair has a strong organismal phenotype. A structurally constrained architecture can still contain redundancy, buffering, and context-dependent function~\cite{Kellis2014}.

%============================================================
\section{Falsifiable Predictions}
\label{sec:predictions}
%============================================================

\begin{prediction}[Mod-8 Intron Non-Uniformity]
\label{pred:mod8}
Preregistered analyses of intron-length distributions should show statistically significant non-uniformity mod~$8$, even though the dominant residue class may depend on transcript weighting or gene set. \emph{Falsifier}: uniform distribution of intron lengths mod $8$ across well-powered datasets.
\end{prediction}

\begin{prediction}[Minimal-Intron Window]
\label{pred:minintron}
The interval $85$--$92$\,bp (centered on $8\ph^5 \approx 88.7$) should be enriched among minimal functional introns in curated mammalian datasets. \emph{Falsifier}: minimal functional introns consistently $\gg 100$ or $\ll 70$\,bp.
\end{prediction}

\begin{prediction}[Topology-Aware Mutation Metrics]
\label{pred:q6}
The \Qsix{} adjacency structure should predict mutational tolerance better than a coarse synonymous/non-synonymous dichotomy, especially for near-neighbor substitutions at the third codon position. \emph{Falsifier}: no improvement in pathogenicity prediction from \Qsix{} distance over existing metrics (BLOSUM, Grantham).
\end{prediction}

\begin{prediction}[Bio-Clocking Clusters]
\label{pred:bioclock}
Biological timing constants involved in folding, transcription, and neural signaling should cluster near $\ph$-spaced rungs on a logarithmic timescale. \emph{Falsifier}: timing constants uniformly or randomly distributed on a log scale.
\end{prediction}

\begin{prediction}[Phase-Sensitive Constraint Beyond Sequence Conservation]
\label{pred:phase}
Regions of non-coding DNA that are not conserved by standard alignment-based methods should nonetheless show measurable constraint when assayed against helical position, chromatin timing, or coordinated expression phasing. \emph{Falsifier}: no detectable phase-sensitive signal in non-conserved non-coding sequence.
\end{prediction}

\begin{prediction}[The Code Optimizes More Than Chemical Similarity --- \textbf{CONFIRMED}]
\label{pred:composite}
\textbf{Original prediction (first draft):} RS predicts the code is locally optimal under a composite $\Jcost$ metric that includes biosynthetic-pathway cost and position-dependent weighting, not just simple chemical distance.

\textbf{Status:} \emph{Substantially confirmed.} Proposition~\ref{thm:global-opt} shows that under a position-weighted cost with 8 components (hydrophobicity, charge, volume, and biosynthetic gap, each weighted by codon position), the standard code is the unique optimum over the 363,004 classwise non-identity permutations currently enumerated. The simple chemistry-only metric fails (42 improving swaps exist~\cite{Washburn2026a}); the position-weighted metric succeeds on the enumerated classwise search with margin $+37{,}136$~\cite{Washburn2026b}.
\end{prediction}

%============================================================
\section{Open Problems}
\label{sec:open}
%============================================================

\paragraph{The codon-assignment problem (now partially resolved).}
At the first draft, the formalism explained why there are 64 codons and 20 amino acids but not why \emph{this} codon maps to \emph{that} amino acid. As of this revision, Proposition~\ref{thm:global-opt} establishes that the standard code is the unique optimum on a 363,004-element classwise degeneracy-preserving search under a position-weighted cost~\cite{Washburn2026b}. The remaining open questions are (i) extending this from the classwise search to the full coupled degeneracy-preserving product space, and (ii) the WToken-to-amino-acid structural bridge: the cardinality match ($20 = 20$) and the fine-class constraint (WToken mode families partition the 20 amino acids into 8 classes of size 2--4) are formalized, but the mechanistic route from DFT-derived semantic classes to amino-acid side-chain chemistry is open.

\paragraph{The chemistry bridge.}
The WToken/amino-acid bijection is a theorem about cardinalities and mode families. The WToken fine-class partition constrains the 20! assignment space to 36,864 alternatives, and the position-weighted cost strongly favors the standard code on the formally and computationally explored neighborhoods~\cite{Washburn2026b}. What remains open is the \emph{physical mechanism} linking DFT mode structure to tRNA recognition, aminoacyl-tRNA synthetase specificity, and biosynthetic pathway organization.

\paragraph{From definitions to derivations.}
The current genome-budget fractions ($1.5/35/25/25/10/3.5$) are explicit definitions in the Lean files, not outputs of a comparative-genomics pipeline or a first-principles optimization. To become a testable theory, they must be replaced by derived quantities.

\paragraph{The intron evidence base.}
The golden-angle theorem and minimal-intron estimate are clean formal results. The mod-8 enrichment claims are explicitly dataset-scoped in the current Lean files; different gene sets yield different peak residues. A publishable biology paper will need preregistered pipelines, raw counts, and sensitivity analyses across multiple annotation releases.

\paragraph{Physical FOLD.}
FOLD is a mathematically clean operator on an 8-register. Whether meiotic recombination, homologous pairing, and chromatin remodeling actually implement it at the molecular level is an open experimental question.

\paragraph{Engagement with transposable-element biology.}
The ledger-balance interpretation of repetitive DNA is the most speculative of the five layers. It must be tested against the mature body of work on selfish DNA~\cite{Orgel1980}, constructive neutral evolution~\cite{Stoltzfus1999}, and the documented regulatory co-option of transposable elements~\cite{Chuong2017}.

%============================================================
\section{Conclusion}
\label{sec:conclusion}
%============================================================

This paper has presented eight machine-checked results and one computation-assisted codon-assignment result that connect the RS algebraic backbone to concrete features of DNA: the codon-space combinatorics ($64 = 8^2$, $14/16$ family-box wobble symmetry, $8$ four-fold degenerate boxes), the wobble-parity non-preservation theorem, the $(8,1,8)$ code optimality, the golden-angle/intron geometry, the bio-clock recurrence, FOLD idempotence with antisymmetric-phase cancellation, the position-specific mutation-vulnerability ordering (wobble safest, position~2 most dangerous), and the position-weighted classwise codon-assignment optimum (Proposition~\ref{thm:global-opt}). It has also proposed a five-layer decomposition of non-coding DNA as a structured research program and stated six falsifiable predictions, one of which (Prediction~\ref{pred:composite}) is now substantially confirmed.

The honest current status is that RS provides a nontrivial formal scaffold---not yet a finished theory of the genome. The exact budget fractions are modeling choices, and the strongest biological interpretations (ledger balance, physical FOLD, phase-sensitive constraint) remain hypotheses. What distinguishes the RS approach from prior accounts of non-coding DNA is not a claim of completeness but a specific algebraic language ($\Jcost$, $\ph$, 8-tick) that generates falsifiable quantitative predictions rather than post-hoc narrative.

The immediate next steps are clear: preregistered intron-length analyses, a \Qsix{}-based mutation-tolerance benchmark, and a comparative study of repeat-element phasing across mammalian clades. Each of these can be conducted independently of the broader RS theory and evaluated on its own empirical merits.

\paragraph{Lean formalization.}
All theorem-level structural results in this paper are machine-checked in Lean~4 (v4.27.0) with Mathlib, within the \texttt{Genetics} namespace of the \texttt{IndisputableMonolith} codebase. The namespace now contains \textbf{49 modules totaling 8,246 lines} with zero \texttt{sorry}, zero \texttt{axiom}, and zero placeholder proofs. The one exception is Proposition~\ref{thm:global-opt}, whose exact classwise search size is formalized in Lean and whose 52 local/classwise certificates are formalized in Lean, while the full 363,004-case optimum statement is confirmed by exhaustive external computation. Key module groups include:
\begin{itemize}[nosep]
\item \textbf{Codon objective and optimization}: \texttt{CodonNeighborCost}, \texttt{MutationModel}, \texttt{CompositeJCost}, \texttt{CodonAssignmentObjective}, \texttt{CodonAssignmentOptimality}, \texttt{CodonAssignmentGlobalOptimality}, \texttt{DegeneracySearchSpace}, \texttt{ComputableAssignmentCost}, \texttt{CertifiedResults}, \texttt{PositionWeightedCost}
\item \textbf{Wobble, parity, and conservation}: \texttt{AminoAcidZClassification}, \texttt{ZInvariantPipeline}, \texttt{ZParityDegeneracyTable}, \texttt{WobbleZParity}, \texttt{ConservationHierarchy}
\item \textbf{Helix geometry and intron placement}: \texttt{IntronStructure}, \texttt{PhaseAlignmentLayer}, \texttt{HelicalPhaseModel}, \texttt{IntronPlacementOptimality}
\item \textbf{Bio-clocking and FOLD}: \texttt{TemporalCoordinationLayer}, \texttt{ChromosomePhaseProfile}, \texttt{HydrationGearbox}, \texttt{GearboxSelectivity}
\item \textbf{Five-layer genome decomposition}: \texttt{ErrorCorrectionLayer}, \texttt{PhaseAlignmentLayer}, \texttt{LedgerBalancingLayer}, \texttt{ModularArchitectureLayer}, \texttt{TemporalCoordinationLayer}, \texttt{GenomeBudget}, \texttt{DNAFramework}
\item \textbf{WToken bridge and search}: \texttt{WTokenAssignmentBridge}, \texttt{BiosyntheticPathways}, \texttt{TranslationKinetics}, \texttt{SearchSpaceCounting}, \texttt{DominanceLemmas}, \texttt{NearUniqueness}
\end{itemize}

\paragraph{Acknowledgments.}
The Lean~4 formalization builds on Mathlib. Empirical intron data referenced by the current genetics files are from Ensembl (GRCh38).

\begin{thebibliography}{99}

\bibitem{Ohno1972}
S.~Ohno, \emph{So much ``junk'' DNA in our genome}, Brookhaven Symp.\ Biol.\ \textbf{23}, 366--370 (1972).

\bibitem{ENCODE2012}
The ENCODE Project Consortium, \emph{An integrated encyclopedia of DNA elements in the human genome}, Nature \textbf{489}, 57--74 (2012).

\bibitem{Graur2013}
D.~Graur \emph{et al.}, \emph{On the immortality of television sets: ``function'' in the human genome according to the evolution-free gospel of ENCODE}, Genome Biol.\ Evol.\ \textbf{5}, 578--590 (2013).

\bibitem{Doolittle2013}
W.F.~Doolittle, \emph{Is junk DNA bunk? A critique of ENCODE}, Proc.\ Natl.\ Acad.\ Sci.\ USA \textbf{110}, 5294--5300 (2013).

\bibitem{Ponting2011}
C.P.~Ponting and R.C.~Hardison, \emph{What fraction of the human genome is functional?}, Genome Res.\ \textbf{21}, 1769--1776 (2011).

\bibitem{Rands2014}
C.M.~Rands \emph{et al.}, \emph{8.2\% of the human genome is constrained: variation in rates of turnover across functional element classes in the human lineage}, PLoS Genet.\ \textbf{10}, e1004525 (2014).

\bibitem{Kellis2014}
M.~Kellis \emph{et al.}, \emph{Defining functional DNA elements in the human genome}, Proc.\ Natl.\ Acad.\ Sci.\ USA \textbf{111}, 6131--6138 (2014).

\bibitem{Orgel1980}
L.E.~Orgel and F.H.C.~Crick, \emph{Selfish DNA: the ultimate parasite}, Nature \textbf{284}, 604--607 (1980).

\bibitem{Stoltzfus1999}
A.~Stoltzfus, \emph{On the possibility of constructive neutral evolution}, J.\ Mol.\ Evol.\ \textbf{49}, 169--181 (1999).

\bibitem{Chuong2017}
E.B.~Chuong, N.C.~Elde, and C.~Feschotte, \emph{Regulatory activities of transposable elements: from conflicts to benefits}, Nat.\ Rev.\ Genet.\ \textbf{18}, 71--86 (2017).

\bibitem{HaigHurst1991}
D.~Haig and L.D.~Hurst, \emph{A quantitative measure of error minimization in the genetic code}, J.\ Mol.\ Evol.\ \textbf{33}, 412--417 (1991).

\bibitem{FreelandHurst1998}
S.J.~Freeland and L.D.~Hurst, \emph{The genetic code is one in a million}, J.\ Mol.\ Evol.\ \textbf{47}, 238--248 (1998).

\bibitem{Horvath2013}
S.~Horvath, \emph{DNA methylation age of human tissues and cell types}, Genome Biol.\ \textbf{14}, R115 (2013).

\bibitem{Washburn2025}
J.~Washburn, \emph{The Algebra of Reality: A Recognition Science Derivation of Physical Law}, Axioms \textbf{15}(2), 90 (2025).

\bibitem{Washburn2026a}
J.~Washburn, \emph{The Standard Genetic Code Is Not Locally Optimal Under Simple Chemical Distance}, in preparation (2026).

\bibitem{Washburn2026b}
J.~Washburn, \emph{The Standard Genetic Code Is Globally Optimal Under Position-Weighted Mutation Cost}, in preparation (2026).

\end{thebibliography}

\end{document}
